% =====================================================================
%  Liga os metadados do latexkit.config.json aos comandos da abntex2.
%
%  As macros \lkAlgumaCoisa vem de config/metadata.tex, que e gerado.
%  \lkIfAlgumaCoisa{ha valor}{esta vazio} omite blocos quando o campo
%  nao foi preenchido.
%
%  Para mudar um valor, edite latexkit.config.json.
% =====================================================================

\titulo{\lkTitle}
\autor{\lkAuthor}
\local{\lkCity}
\data{\lkYear}
\tipotrabalho{\lkReportType}

% A NBR 10719 pede a identificacao da instituicao responsavel e, quando
% houver, da unidade e do cliente ou orgao demandante.
\instituicao{%
  \lkInstitution
  \lkIfUnit{\par\lkUnit}{}%
}

% Preambulo da folha de rosto: a que o relatorio se destina.
\preambulo{%
  \lkReportType%
  \lkIfProject{ referente ao projeto \lkProject}{}%
  \lkIfClient{, elaborado para \lkClient}{}%
  .%
  \lkIfPeriod{\par\vspace{0.5em}Periodo de referencia: \lkPeriod.}{}%
}

% Em um relatorio, o campo "orientador" da abntex2 vira o responsavel
% tecnico. A folha de rosto imprime o rotulo sem checar se ha nome, entao
% o campo vazio precisa zerar tambem o rotulo — do contrario sobraria um
% "Responsavel tecnico:" solto na pagina.
\lkIfSupervisor{\orientador[Responsavel tecnico:]{\lkSupervisor}}{\orientador[]{}}

\makeatletter
\hypersetup{
  pdftitle={\@title},
  pdfauthor={\@author},
  pdfsubject={\lkReportType},
  pdfkeywords={\lkKeywords},
  pdfcreator={latexkit}
}
\makeatother
